\chapter{The Parrot Who Wouldn't Talk}
'You're no beauty! Can't talk, can't sing, can't dance!' With these words, Aunt Ruby would taunt the unfortunate parakeet who glared morosely at the everyone form his ornamental cage at one end of the long veranda of Granny's bungalow in north Indian.
\par
In those distant days, almost everyone-Indian or European-kept a pet parrot or parakeet, or 'lovebird' as some of the smaller ones were called. Sometimes there birds became great talkers, or mantras (religious chants) or admonitions to the children of the house, such as 'Padho, beta, padho!' ('Study, child, study!') or, for the benefit of boys like me-'Don't be greedy, don't be greedy!'
\par
The expression were, of course, picked up by the parrot over a period of time, after many repetitions by whichever member of the household had taken on the task of teaching the bird to talk.
\par
But our parrot refused to talk.
\par
He'd been bought by Aunt Ruby form a bird catcher who'd visited all the house on our road, selling caged birds ranging from colourful budgerigars to chirpy little munias and even common sparrows that had been dabbed with paint and passed off as some exotic species. Neither Granny nor Grandfather were keen con keeping caged birds as pets, but Aunt Ruby threatened to throw a tantrum if she did not get her way-and Aunt Ruby's tantrums were dreadful to behold!
\par
Anyway, she insisted on keeping the parrot and teaching it to talk. But the bird took an instant dislike to my aunt and resisted all her blandishments.
\par
'Kiss, kiss!' Aunt Ruby would coo, putting her face close to the barge of the cage. But the parrot would back away, its beady little eyes getting even smaller with anger at the prospect of being kissed by Aunt Ruby. On one occasion, it lunged forward without warning and knocked my aunt's spectacles off her nose.
\par
After that, Aunt Ruby gave up her endearments and became quiet hostile towards the poor bird, making faces at it and calling out, 'Can't talk, can't sing, can't dance!' and other nasty comments.
\par
It fell upon me, then ten years old, to feed the parrot, and it seemed quite happy to receive green chillies and ripe tomatoes from my hands, these delicacies being supplemented by slices of mango, for it was then the mango season. It also gave me an opportunity to consume a couple of mangoes while feeding the parrot.
\par
One afternoon, while everyone was indoors enjoying a siesta, I gave the parrot his lunch and the deliberately left the cage door open. Seconds later, the bird was winging its way to the freedom of the mango orchard.
\par
At the same time, Grandfather came to the veranda and remarked, 'I see your aunt's parrot has escaped!'
\par
'The door was quite loose,' I said with a shrug. 'Well, I don't suppose we'll see it again.'
\par
Aunt Ruby was upset at first and threatened to buy another bird. We put her off by promising to buy her a bowl of goldfish.
\par
'But goldfish don't talk!' she protested.
\par
'Well, neither did you bird,' said Grandfather. 'So we'll get you a gramophone. You can listen to Clara Cluck all day. They say she sings like a nightingale.'
\par
I thought we'd never see the parrot again, but it probably missed its green chillies, because a few days later I found the bird sitting on the veranda railing, looking expectantly at me with its head cocked to one side, Unselfishly, I gave the parrot half of my mango.
\par
While the bird was enjoying the mango, Aunt Ruby emerged from her room and, with a cry of surprise, called out, 'Look, my parrot's came back! He must have missed me!'
\par
With a loud squawk, the parrot flew out of her reach and, parching on the nearest rose bush, glared at Aunt Ruby and shrieked at her in my aunt's familiar tones: You're no beauty!cCan't talk, can't sing, can't dance!'
\par
Aunt Ruby went ruby-red and dashed indoors.
\par
But that wasn't the end of the affair. The parrot became a frequent visitor to the garden and veranda and whenever it saw Aunt Ruby it would call out, 'You're no beauty, you're no beauty! Can't talk, can't sing, can't dance!'
\par
The parrot had learnt to talk, after all.




